\section{Clarification of Notions: Feature, Representation, Concept, and Semantics}

In the deep learning literature, the notions of \textit{feature}, \textit{representation}, \textit{concept}, and \textit{semantics} are often used with varying degrees of precision across communities (machine learning, computer vision, and interpretability).
This section provides a conceptual clarification grounded in foundational works in representation learning.

\subsection{Feature}

A \textbf{feature} is a numerical variable derived from observed data, constituting the elementary unit of information processed by a learning model.
In classical machine learning, features were manually engineered.
In deep learning, they are now automatically learned at multiple levels of abstraction within the network \cite{lecun2015deep}.
Formally, a feature can be interpreted as a component of an input vector $x \in \mathbb{R}^d$ or of an intermediate activation $h^{(l)} \in \mathbb{R}^{d_l}$.
Thus, features correspond to local, scalar-valued quantities forming the basis of learned representations.

\subsection{Representation}

A \textbf{representation} refers to a transformed encoding of data in a latent space, obtained through a composition of parameterized functions:
\[
h^{(l)} = f^{(l)}(h^{(l-1)})
\]
Each layer of a deep neural network learns such a transformation, producing a hierarchy of representations of increasing abstraction.
According to Bengio et al.~\cite{bengio2013representation}, deep learning aims at learning distributed representations that capture explanatory factors of variation in the data.

These representations are typically:
\begin{itemize}
    \item distributed (each dimension contributes to multiple factors),
    \item hierarchical,
    \item task-dependent.
\end{itemize}

From a geometric perspective, a representation corresponds to a learned embedding of the data into a latent space where relevant factors become more explicitly organized.

\subsection{Concept}

A \textbf{concept} is defined as a semantic abstraction corresponding to a stable regularity in the real world, independent of any particular encoding within a model.
In neural networks, concepts are not explicitly represented but rather emerge from the structure of learned representations.
Work in interpretability has shown that concepts can often be associated with directions or subspaces in latent representations \cite{bau2017network, kim2018tcav}.
For instance, a concept such as \textit{``dog''} may correspond to multiple internal configurations without a unique or bijective mapping to a specific neuron or feature.
Thus, concepts are best understood as emergent semantic entities inferred from learned representations rather than explicitly encoded variables.

\subsection{Semantics}

\textbf{Semantics} refers to the relationship between an internal representation and its interpretation in terms of meaning in the real world.
In deep learning, semantics are typically implicit and must be inferred a posteriori.
They correspond to the alignment between geometric structures in representation spaces and human-interpretable concepts.
Representation learning has shown that certain geometric properties (e.g., clustering structure, linear separability, vector directions) correlate with semantic regularities \cite{mikolov2013distributed, goodfellow2016deep}.